\documentclass[tikz,border=3pt]{standalone}
\usetikzlibrary{arrows.meta}
\begin{document}
\begin{tikzpicture}[>=Latex, font=\small]
  \coordinate (O) at (0,0);
  \coordinate (F) at (1.3,0);   % pie de la perpendicular = punta de c w1
  \coordinate (W1) at (4.2,0);  % punta de w1 = v1
  \coordinate (W2) at (1.3,2.2); % punta de w2

  \draw[-{Latex[length=2.5mm]}] (O) -- (W1) node[right] {$w_1=v_1$};
  \draw[-{Latex[length=2.5mm]}] (O) -- (W2) node[above left] {$w_2$};
  \draw[-{Latex[length=2.5mm]}] (F) -- (W2) node[midway, left] {$v_2$};
  \draw (W2) -- (W1);

  \draw (F) ++(-0.18,0) -- ++(0,0.18) -- ++(0.18,0);

  \node[below] at (F) {$cw_1$};
\end{tikzpicture}
\end{document}
